\chapter{Lecture 4: Shell Scripting, Testing and Intro to C++}
\begin{center}
\textit{January 16, 2018}
\end{center}

\section{Conditions with \code{[ ]} (continued)}
Recall that \code{[ cond ]} (the \code{test} program) evaluates a condition and
sets its status code to 0 (true) or 1 (false). Output we do not care about is
sent to \code{/dev/null}.

\begin{center}
\begin{tabular}{@{}lp{0.62\textwidth}@{}}
\toprule
Test & True when \\
\midrule
\code{[ s1 = s2 ]}, \code{[ s1 == s2 ]} & strings are equal \\
\code{[ s1 != s2 ]} & strings are different \\
\code{[ -z s ]}, \code{[ -n s ]} & string is empty / non-empty \\
\code{[ a -eq b ]}, \code{-ne} & integers equal / not equal \\
\code{-gt}, \code{-ge}, \code{-lt}, \code{-le} & integer $>$, $\ge$, $<$, $\le$ \\
\code{[ -e f ]} & file \code{f} exists \\
\code{[ -f f ]}, \code{[ -d f ]} & \code{f} is an ordinary file / a directory \\
\code{[ -r f ]}, \code{-w}, \code{-x} & \code{f} is readable / writable / executable \\
\code{[ c1 -a c2 ]}, \code{[ c1 -o c2 ]} & and / or \\
\code{[ ! c ]} & not \\
\bottomrule
\end{tabular}
\end{center}

\begin{important}[Strings vs.\ numbers]
\code{=} and \code{!=} compare \emph{strings}; \code{-eq}, \code{-lt}, \dots{}
compare \emph{integers}. \code{[ 01 = 1 ]} is false but \code{[ 01 -eq 1 ]} is
true. Always quote variables inside tests (\code{[ "\$\{x\}" = "abc" ]}) so an
empty variable doesn't turn into a syntax error.
\end{important}

\section{If Statements}
\begin{lstlisting}[style=bash]
if [ cond ]; then
    statements
elif [ cond ]; then
    statements
else
    statements
fi
\end{lstlisting}
The thing after \code{if} is really any command: the branch is taken if its
status code is 0. So we can test a command directly, without looking at
\code{\$?}.

\begin{example}
Good-password checker, with output, and a usage message if called incorrectly.
\begin{lstlisting}[style=bash]
#!/bin/bash
usage() {
    echo "Usage: $0 password" 1>&2     # error messages go to stderr
    exit 1
}

if [ $# -ne 1 ]; then
    usage
fi

if egrep "^$1$" /usr/share/dict/words > /dev/null; then
    echo "Not a good password"
else
    echo "Maybe a good password"
fi
\end{lstlisting}
Notes:
\begin{itemize}
  \item \code{\$\#} is the number of arguments.
  \item \code{usage() \{ ... \}} defines a shell \emph{function}; inside a
        function, \code{\$1}, \code{\$2}, \dots{} refer to the function's own
        arguments.
  \item \code{exit 1} ends the script with a non-zero (failure) status.
  \item \code{1>\&2} sends the message to stderr.
\end{itemize}
\end{example}

\section{Loops}

\subsection{While Loops}
\begin{lstlisting}[style=bash]
while [ cond ]; do
    statements
done
\end{lstlisting}

\begin{example}
Print the numbers from 1 to \code{\$1}.
\begin{lstlisting}[style=bash]
#!/bin/bash
x=1
while [ ${x} -le $1 ]; do
    echo ${x}
    x=$((x + 1))
done
\end{lstlisting}
Variables hold strings, so \code{x=x+1} would just store the string
\code{x+1}. Arithmetic is done with \code{\$((...))}.
\end{example}

\subsection{For Loops}
A \code{for} loop iterates over a list of words:
\begin{lstlisting}[style=bash]
for x in a b c d; do
    echo ${x}
done
\end{lstlisting}

\begin{example}
Rename every \code{.cpp} file in the current directory to have a \code{.cc}
extension (the task the graphical shell couldn't do in Lecture~1).
\begin{lstlisting}[style=bash]
#!/bin/bash
for name in *.cpp; do
    mv "${name}" "${name%cpp}cc"
done
\end{lstlisting}
The glob \code{*.cpp} expands to the list of matching files. \code{\$\{name\%cpp\}}
is the value of \code{name} with the trailing \code{cpp} removed.
\end{example}

\begin{example}
How many times does the word \code{\$1} appear in the file \code{\$2}?
\begin{lstlisting}[style=bash]
#!/bin/bash
x=0
for word in $(cat "$2"); do
    if [ "${word}" = "$1" ]; then
        x=$((x + 1))
    fi
done
echo ${x}
\end{lstlisting}
\code{\$(cat "\$2")} is deliberately unquoted so that it splits into one word
per loop iteration.
\end{example}

\begin{example}
Payday is the last Friday of the month. On what day is this month's payday?
\begin{lstlisting}[style=bash]
#!/bin/bash
answer() {
    if [ "$1" -eq 31 ]; then
        echo "This month: the 31st"
    else
        echo "This month: the ${1}th"
    fi
}
# cal prints a calendar; column 6 is Friday. Take the last non-empty one.
answer $(cal | awk '{print $6}' | egrep "[0-9]" | tail -1)
\end{lstlisting}
Generalizing to any month/year: pass \code{\$1 \$2} along to \code{cal}
(\code{cal \$1 \$2}); with no arguments \code{cal} uses the current month.
\end{example}

\section{Testing}
Testing is an essential part of program development. It is \emph{not}
debugging: testing reveals the presence of bugs; it cannot prove their
absence. Testing should happen \emph{while} you write code, not only at the end.

\begin{itemize}
  \item \textbf{Human testing}: code inspections, walkthroughs, explaining your
        code to someone else.
  \item \textbf{Machine testing}: running the program on inputs and checking the
        outputs.
  \item \textbf{Black-box testing}: tests are written from the specification
        only, without looking at the implementation. Cover various classes of
        input (numeric ranges, positive/negative), boundary cases (edges of
        ranges, empty input, first/last), multiple simultaneous boundaries
        (corner cases), intuition, and things that are clearly wrong (invalid
        input).
  \item \textbf{White-box testing}: tests are written with the code in view, so
        that every branch, loop and function is exercised (statement and branch
        coverage).
  \item \textbf{Performance testing}: is the program efficient enough?
  \item \textbf{Regression testing}: keep a test suite and re-run it every time
        the code changes, to make sure nothing that used to work is now broken.
\end{itemize}

\begin{center}
\begin{tikzpicture}[>=Stealth,font=\small]
  \node[draw,fill=black!80,text=white,minimum width=2cm,minimum height=1.2cm] (bb) at (0,0) {program};
  \draw[->] (-2.2,0.3) -- node[above,font=\scriptsize]{input} (-1,0.3);
  \draw[->] (1,0.3) -- node[above,font=\scriptsize]{output} (2.2,0.3);
  \node[font=\scriptsize,align=center] at (0,-1.1) {black box: tests come\\ from the specification};
  \begin{scope}[xshift=6.5cm]
  \draw (-1.3,-0.8) rectangle (1.3,0.8);
  \node[draw,diamond,aspect=1.6,inner sep=1pt,font=\scriptsize] (d) at (0,0.25) {$?$};
  \draw[->,blue!60!black,thick] (d.west) -- ++(-0.6,-0.6);
  \draw[->,red!60!black,thick] (d.east) -- ++(0.6,-0.6);
  \draw[->] (-2.2,0.25) -- node[above,font=\scriptsize]{input} (-1.3,0.25);
  \draw (-1.3,0.25) -- (d.west);
  \node[font=\scriptsize,align=center] at (0,-1.3) {white box: tests chosen so\\ every branch is taken};
  \end{scope}
\end{tikzpicture}
\end{center}

A1 asks us to automate regression testing with a script. A minimal version: for
each test stem listed in a suite file, run the program on \code{stem.in} and
compare with \code{stem.out}.
\begin{lstlisting}[style=bash]
#!/bin/bash
# usage: ./runSuite suite-file program
for stem in $(cat "$1"); do
    tmp=$(mktemp)
    "$2" < "${stem}.in" > "${tmp}"
    if ! diff "${tmp}" "${stem}.out" > /dev/null; then
        echo "Test failed: ${stem}"
    fi
    rm "${tmp}"
done
\end{lstlisting}
For each test stem the script does the following:
\begin{center}
\begin{tikzpicture}[>=Stealth,font=\small,node distance=1.1cm,
  file/.style={draw,fill=yellow!10,minimum height=0.7cm},
  prog/.style={draw,fill=gray!15,minimum height=0.8cm,rounded corners}]
  \node[file] (in) {\code{stem.in}};
  \node[prog,right=of in] (p) {program \code{\$2}};
  \node[file,right=of p] (tmp) {\code{\$\{tmp\}}};
  \node[prog,right=of tmp] (d) {\code{diff}};
  \node[file,below=0.5cm of d] (out) {\code{stem.out}};
  \node[right=0.7cm of d,align=left,font=\scriptsize] (r) {differ $\Rightarrow$\\ ``Test failed''};
  \draw[->] (in) -- node[above,font=\scriptsize]{stdin} (p);
  \draw[->] (p) -- node[above,font=\scriptsize]{stdout} (tmp);
  \draw[->] (tmp) -- (d); \draw[->] (out) -- (d); \draw[->] (d) -- (r);
\end{tikzpicture}
\end{center}

\section{Introduction to C++}
C++ was designed by Bjarne Stroustrup (early 1980s) as ``C with classes''. It is
(almost) a superset of C, adding classes, references, operator and function
overloading, templates, exceptions, a much bigger standard library, and more.
In CS 246 we use the C++14 standard.

\begin{lstlisting}[style=cpp]
#include <iostream>
using namespace std;

int main() {
    cout << "Hello world" << endl;
    return 0;
}
\end{lstlisting}

Differences from C that are already visible:
\begin{itemize}
  \item \code{\#include <iostream>} instead of \code{<stdio.h>}; I/O is done with
        \emph{streams} (\code{cout}) and the operator \code{<<} instead of
        \code{printf}.
  \item \code{using namespace std;} lets us write \code{cout} instead of
        \code{std::cout}.
  \item \code{main} must return \code{int}; \code{return 0} may be omitted in
        \code{main} (0 is returned by default).
  \item \code{endl} writes a newline and flushes the output buffer.
\end{itemize}

Compiling and running:
\begin{lstlisting}[style=bash]
g++ -std=c++14 hello.cc            # produces a.out
./a.out
g++ -std=c++14 hello.cc -o hello   # name the executable
./hello
\end{lstlisting}
(On the student environment \code{g++14} is an alias for
\code{g++ -std=c++14}.) C++ source files conventionally use \code{.cc} (or
\code{.cpp}); headers use \code{.h}.

\begin{aside}
In CS 246 we write C++, not C: use \code{new}/\code{delete} rather than
\code{malloc}/\code{free}, and C++ streams and strings rather than
\code{printf}/\code{scanf} and \code{char*}. Details start next lecture.
\end{aside}

\section*{Summary}
\begin{itemize}
  \item \code{[ ]} tests strings (\code{=}, \code{!=}), integers (\code{-eq},
        \code{-lt}, \dots) and files (\code{-e}, \code{-f}, \code{-d}).
  \item \code{if/elif/else/fi}, \code{while ... do ... done},
        \code{for x in list; do ... done}; arithmetic with \code{\$((...))};
        functions and \code{exit}.
  \item Black-box, white-box and regression testing; automate with scripts.
  \item C++ hello world, \code{iostream}, and compiling with
        \code{g++ -std=c++14}.
\end{itemize}
